\section{Complete results and fitting budgets}
\label{app:allresults}
\subsection{Individual models, baselines, and ensembles}
\label{app:matchedbaselines}
Table~\ref{tab:allexperts} reports all nine library members and the three
ensemble rules on every dataset. Table~\ref{tab:baseline_detail}
completes that comparison with individual
additional baseline errors. The latter include the parameter kNN and training
mean controls, which are outside the ranked B1 to B11 collection. Every ranked
method covers all 22 datasets. Errors average the same evaluation cases
within each partition, followed by the three partitions. ERA5 uses its separate
seven case evaluation set. Full precision values, per case errors, and replay
metadata are retained in the companion package.

\begin{table}[h!]
\centering\footnotesize
\caption{\textbf{Additional baseline errors across all 22 datasets.} Mean relative $L^2$
errors in percent. B1: Autoregressive POD GP. B2: Nonlinear POD GP.
B3: Composite MF MLP. B4: Composite MF DeepONet. B5: Latent MF network.
B6: Nonlinear decoder transfer. B7: MFRNP adapter. B8: Fidelity basis FNO.
B9 and B10: FNO transfer I and II. B11: Affine POD GP. B12: Parameter kNN.
B13: Training mean. Mixtures use five fitting examples and the same evaluation
cases. A dash marks an unavailable additional
control, not a missing ranked method. B9 and B10's implementation relationship
is described in Appendix~\ref{app:adaptations}. Imported data are from
MFRNP \citep{niu2024} and ERA5 reanalysis \citep{hersbach2020}.
Daggers identify the data quality exclusions in Appendix~\ref{app:subsets}.}
\label{tab:baseline_detail}
{\scriptsize\setlength{\tabcolsep}{2pt}
\renewcommand{\arraystretch}{1.15}
\resizebox{\linewidth}{!}{\begin{tabular}{lrrrrrrrrrrrrrrr}
\toprule
Dataset & B1 & B2 & B3 & B4 & B5 & B6 & B7 & B8 & B9 & B10 & B11 & B12 & B13 & \shortstack[r]{Inverse\\error\\mixture} & \shortstack[r]{Fitted\\mixture} \\
\midrule
Darcy & 3.37 & 3.66 & 4.03 & 6.66 & 5.91 & 5.2 & 12.4 & 5.55 & 5.27 & 5.28 & 3.6 & 21.5 & 38 & 2.55 & 2.33 \\
Eikonal & 3.49 & 2.99 & 3.12 & 3.35 & 3.38 & 2.65 & 3.62 & 2.21 & 2.33 & 2.34 & 4.88 & 5.51 & 31.8 & 1.76 & 1.76 \\
Helmholtz & 12.3 & 15.2 & 11.5 & 7.84 & 8.12 & 99.6 & 58.6 & 99.7 & 4.12 & 4.12 & 15.9 & 25.7 & 92 & 6.89 & 6.55 \\
Poisson I & 0.00731 & 0.00531 & 0.181 & 0.161 & 0.1 & 0.479 & 6.79 & 0.239 & 0.0804 & 0.0806 & 2.68 & 7.51 & 31.9 & 0.008 & 0.00788 \\
Poisson II$^{\dagger}$ & 0.0128 & 0.0306 & 0.635 & 1.72 & 0.849 & 3.93 & 11.7 & 82.1 & 0.355 & 0.357 & 2.01 & 12.3 & 38.5 & 0.0125 & 0.0129 \\
Pressure Poisson$^{\dagger}$ & 0.457 & 0.472 & 0.07 & 1.53 & 0.127 & 0.446 & 1.05 & 0.684 & 0.0639 & 0.0645 & 0.681 & 0.196 & 1.58 & 0.0769 & 0.0802 \\
Heat I & 0.326 & 0.327 & 0.0696 & 0.125 & 0.069 & 0.432 & 1.29 & 0.275 & 0.0207 & 0.0206 & 1.09 & 1.35 & 12.4 & 0.00593 & 0.0062 \\
Heat II & 0.312 & 0.312 & 0.111 & 0.183 & 0.0721 & 0.219 & 7.4 & 0.32 & 0.0202 & 0.0201 & 0.405 & 0.856 & 13.2 & 0.00437 & 0.00435 \\
Cavity & 0.307 & 0.409 & 23.9 & 4.99 & 0.0419 & 24.1 & 55.7 & 0.129 & 0.622 & 0.666 & 37.2 & 0.077 & 17.1 & 0.142 & 0.142 \\
Navier Stokes & 2.8 & 2.48 & 2.08 & 4.01 & 3.71 & 7.08 & 25.7 & 1.62 & 1.52 & 1.52 & 21.6 & 3.27 & 41.2 & 1.52 & 1.5 \\
Rayleigh Bénard & 0.283 & 0.295 & 0.0542 & 0.0588 & 0.0134 & 0.0613 & 4.5 & 0.78 & 0.0113 & 0.011 & 0.715 & 0.0575 & 9.1 & 0.00286 & 0.00269 \\
Burgers & 2.16 & 2.5 & 1.54 & 8.14 & 1.22 & 15.7 & 27.7 & 0.928 & 0.562 & 0.562 & 11.7 & 0.677 & 73.2 & 0.513 & 0.502 \\
Euler & 5.6 & 5.54 & 6.56 & 6.16 & 6.15 & 8.81 & 9.15 & 7 & 6.52 & 6.43 & 6.25 & 6.21 & 32.2 & 5.64 & 5.73 \\
Porous medium & 1.15 & 0.859 & 0.299 & 1.31 & 0.772 & 0.918 & 33.7 & 5.14 & 0.0718 & 0.0718 & 1.98 & 1.92 & 31.9 & 0.0262 & 0.0263 \\
Allen Cahn & 55.5 & 98.6 & 59.9 & 44.8 & 41.2 & 94.4 & 69 & 48.9 & 55.8 & 55.8 & 55.4 & 78.3 & 100 & 37.9 & 34.6 \\
Cahn Hilliard I$^{\dagger}$ & 29.7 & 27.8 & 53.5 & 48.7 & 27.6 & 61.1 & 74.7 & 14.2 & 13.4 & 13.3 & 72.3 & 20.6 & 99.9 & 11.3 & 10.7 \\
Cahn Hilliard II & 64.8 & 83.5 & 85 & 53.7 & 49.8 & 72.4 & 66.7 & 54.4 & 54 & 54 & 64.8 & 78 & 100 & 47.2 & 49.3 \\
Fisher KPP & 6.86 & 6.86 & 6.48 & 4.01 & 3.45 & 7.44 & 4.08 & 4.74 & 3.24 & 3.24 & 6.86 & 5.67 & 6.86 & 2.53 & 2.3 \\
Phase field crystal & 83.5 & 86.4 & 90.7 & 87.5 & 87.8 & 152 & 89.4 & 102 & 86.8 & 86.8 & 83.5 & 86.7 & 87.5 & 83.8 & 83.3 \\
Shallow water & 1.14 & 1.02 & 0.467 & 0.878 & 0.691 & 0.883 & 5.59 & 0.806 & 0.403 & 0.403 & 2.94 & 1.22 & 25.1 & 0.357 & 0.389 \\
Wave & 27.7 & 24.9 & 3.3 & 42.3 & 46.6 & 5.39 & 93.9 & 25.9 & 6.6 & 6.6 & 50.3 & 42.8 & 98.7 & 4.39 & 4.27 \\
\midrule ERA5 & 5 & 6.74 & 45 & 16.8 & 5.05 & 23.9 & 19.5 & 5.15 & 11.9 & 11.9 & 8.45 & 5.16 & 6.74 & 5.81 & 6.49 \\
\bottomrule
\end{tabular}
}}
\end{table}

\subsection{Library quality and the benefit of averaging}
\label{app:headlines}
Two comparisons separate the quality of the available library from the
benefit of combining it. The best library reference minimizes partition
averaged error over M1 to M9 separately for each dataset. The best additional
baseline minimizes the corresponding error over B1 to B11. Both use evaluation
answers retrospectively and are stronger references than the deployable
Selected model, which uses only fitting examples.
Table~\ref{tab:headlinecomparisons} reports these comparisons on the 21 PDE
datasets. Table~\ref{tab:paperbaselines} additionally includes ERA5 in its
dataset ratios and Elo ranking.

\begin{table}[h!]
\centering\small
\caption{\textbf{Library quality and mixture gains on the 21 PDE datasets.}
Ratios are geometric means. Wins and losses require a relative difference
greater than 1\%. Best library and best additional baseline refer to the
retrospective per dataset minima described above.}
\label{tab:headlinecomparisons}\begin{tabular}{lrrrr}
\toprule
Comparison & Class ratio & Dataset ratio & Wins & Losses \\
\midrule
Best library / best additional baseline & 0.766 & 0.834 & 15 & 5 \\
Fitted mixture / best library & 0.923 & 0.989 & 13 & 6 \\
Inverse error mixture / best library & 0.927 & 1.002 & 12 & 7 \\
Fitted mixture / best additional baseline & 0.706 & 0.825 & 13 & 7 \\
Fitted mixture / Selected model & 0.884 & 0.923 & 15 & 5 \\
\bottomrule
\end{tabular}

\end{table}

A still stronger control chooses a potentially different best model in each
partition using its evaluation answers. Against this reference, the fitted
mixture at $K=5$ has class ratio 0.933 and dataset ratio 1.006, with
13 improvements and 6 losses greater than 1\%. Thus its advantage over
retrospective model choice depends on both the reference and aggregation.

\subsection{Fitting budgets and uniform averaging}
\label{app:budgets}
Table~\ref{tab:main} compares the main rules with uniform averaging on the
21 PDE datasets. Table~\ref{tab:budgets} varies the fitting budget while keeping
the models and case partition protocol fixed. Each ratio compares methods that
receive the same number of fitting examples. These summaries measure sensitivity
to ensemble fitting, not independent training seed variation.

\begin{table}[h!]
\centering\small
\caption{\textbf{Weighting the same nine model library.} Ratios compare with
the Selected model at $K=5$ on 21 PDE datasets. Lower is better. Wins and losses
require changes larger than 1\%.}
\label{tab:main}\begin{tabular}{lrrrr}
\toprule
Method & Class balanced & Equal dataset & Wins & Losses \\
\midrule
Selected model & 1.000 & 1.000 & 0 & 0 \\
All pairs FNO & 1.333 & 1.502 & 3 & 15 \\
Uniform (all nine) & 5.919 & 4.944 & 2 & 19 \\
Inverse error mixture & 0.888 & 0.935 & 13 & 5 \\
Fitted mixture & 0.884 & 0.923 & 15 & 5 \\
\bottomrule
\end{tabular}

\end{table}

\begin{table}[h!]
\centering\small
\caption{\textbf{Increasing the fitting budget on the 21 PDE datasets.}
Class aggregated error ratios. The models remain fixed, and every comparison
uses the same allowance of $K$ additional fine fields.}
\label{tab:budgets}\begin{tabular}{rrrr}
\toprule
$K$ & \shortstack[r]{Inverse error mixture /\\Selected model} & \shortstack[r]{Fitted mixture /\\Selected model} & \shortstack[r]{Fitted mixture /\\Inverse error mixture} \\
\midrule
5 & 0.8876 & 0.8840 & 0.9958 \\
10 & 0.8921 & 0.8765 & 0.9826 \\
20 & 0.8900 & 0.8692 & 0.9766 \\
\bottomrule
\end{tabular}

\end{table}

\subsection{Elo ranking}
\label{app:elo}
Each pair of methods plays one comparison per dataset using its unrounded
partition averaged error. Lower error wins and an absolute difference below
$10^{-12}$ is a tie. The nine library members, eleven baseline entries, and
three rules all enter on the same 22 datasets. Ratings start at 1500. For
ratings $R_a,R_b$, the expected score and update are
\begin{equation}
 p_a=\frac{1}{1+10^{(R_b-R_a)/400}},\qquad
 R_a'=R_a+32(s_a-p_a),\qquad R_b'=R_b-32(s_a-p_a).
\end{equation}
Here $s_a$ is one for a win, zero for a loss, and one half for a tie.
Primes denote updated ratings. We reset ratings and independently shuffle
comparisons 500 times, then report the mean final rating. The code fixes the
random orders. Sorting uses unrounded ratings and display uses integers.
Datasets receive equal weight. Elo summarizes how often a method wins rather
than the magnitude of improvement, statistical significance, or training cost.
